%%%%%%%%%%%%%%%%%%%%%%%%%%%%%%%%%%%%%%%%%%%%%%%%%%%%%%%%%%%%%%%%%%%%%%%%%%%
% Englische Fassung des Uebungsblatts zu Lehreinheit 1 (Teil A und B).
% Vorlage: handouts/LE01.tex. Nur in handouts.tex mit \def\wuzsprache{en}.
%%%%%%%%%%%%%%%%%%%%%%%%%%%%%%%%%%%%%%%%%%%%%%%%%%%%%%%%%%%%%%%%%%%%%%%%%%%
\begin{uebungsblatt}{Wonder, Explanation, Method}

\aufgabe{Sorting hypotheses}[in class, small groups, 20 min.]
Copy the explanations from the board into the table. Assign each one to a
stage and, for each explanation, name an observation that would
\emph{refute} it. The three stages: \emph{wonder} (your own model does not
fit), \emph{explanation} (a hypothesis, not yet evaluated), \emph{method}
(a procedure that decides between hypotheses).

\begin{center}
\renewcommand{\arraystretch}{2.2}
\begin{tabular}{|p{0.38\textwidth}|p{0.14\textwidth}|p{0.38\textwidth}|}
\hline
\textbf{Statement / explanation} & \textbf{Stage} & \textbf{What would refute it?} \\
\hline
 & & \\ \hline
 & & \\ \hline
 & & \\ \hline
 & & \\ \hline
\end{tabular}
\end{center}

\aufgabe{Observing two conditions}[in class, individual work]
The experiment is repeated in class under two changed conditions.
\begin{enumerate}
  \item For each condition, note what was changed and what you observed.
  \item Which of your explanations from Task~1 does this support, and which
    does it weaken?
\end{enumerate}
\schreiblinien{4}
\end{uebungsblatt}

\begin{uebungsblattb}{3}{Belief, Alief and Your Own Doubt}

\aufgabe{Belief or alief?}[homework]
A \emph{belief} is a conscious conviction that can be changed by
arguments; an \emph{alief} is an automatic, bodily readiness to react that
keeps running alongside it. Give two examples from your own life in which
what you know and what you experience come apart (example: you know that a
glass bridge will hold you, and still you do not dare to step onto it).
Which procedure could help you follow your knowledge rather than your
experience?
\schreiblinien{4}

\aufgabe{Short essay}[homework, 500 words, due by Unit~2]
\enquote{Describe an experience that made you doubt a conviction for the
first time. What does this have to do with science?} Use at least two of
the key terms of this unit (prediction error, predictive processing,
belief, alief, ideomotor effect).

\aufgabe{Preparing for Unit~2}[homework]
Bring a question that \emph{really} occupies you -- from your studies,
work or everyday life. In Unit~2 it will be sharpened into a research
question.

\begin{lehrhinweis}[Expected answers]
Task~1: at least one explanation per stage; refutation conditions must be
observable (\enquote{the pendulum also swings with the forearm resting on
the table} rather than \enquote{it is not true}). Task~2: the conditions
are closing the eyes and resting the forearm. Eyes closed $\to$ no visual
feedback, the swing decreases; arm resting $\to$ the muscle impulses are
absorbed. Both point to the hand as the cause and weaken every explanation
that assumes a force outside the person. Task~3: examples with fear of
heights, placebo, optical illusions; procedure = measurement, protocol,
blind condition.
\end{lehrhinweis}
\end{uebungsblattb}
